\chapter{Introduzione}
\label{cap:introduzione}

\section{Obiettivo del progetto}

Questa relazione documenta le attivit\`a di \textbf{manutenzione preventiva}
svolte su un esacottero \textbf{DJI F550} equipaggiato con \textbf{Pixhawk 6X}
e firmware \textbf{PX4}, nell'ambito dell'omonimo corso (A.A. 2025/2026).
Il progetto non si limita alla manutenzione in senso stretto: esso ripercorre
l'intero ciclo di \emph{messa in funzione} del velivolo, dalla configurazione
del flight controller fino all'esecuzione di una campagna di voli mirata
all'\emph{acquisizione di dati} utili a caratterizzare lo stato di salute del
sistema.

Il filo conduttore del lavoro \`e riassumibile in tre fasi:
\begin{enumerate}
  \item \textbf{mettere in funzione} il drone, partendo da una piattaforma
        hardware da assemblare e diagnosticare;
  \item \textbf{configurarlo} tramite QGroundControl e i parametri PX4, fino a
        ottenere un velivolo volabile e strumentato per il logging;
  \item \textbf{eseguire i test acquisendo dati}, conducendo una campagna di
        voli che comprende baseline a pale sane e prove con pale volutamente
        danneggiate, su traiettorie di hover e quadrato.
\end{enumerate}

La prospettiva della manutenzione preventiva, tema del corso, attraversa tutte e
tre le fasi: l'attenzione \`e rivolta a \emph{strumentare} il velivolo --- cio\`e
a predisporre l'acquisizione delle grandezze che permettono di riconoscere per
tempo una degradazione --- pi\`u che a riparare a guasto avvenuto. In quest'ottica
ogni anomalia incontrata (dropout GPS, blocco del flight controller, incidente in
volo) viene documentata insieme alla diagnosi e all'azione correttiva adottata.

\section{Contesto del corso e metodo di lavoro}

Il lavoro \`e stato condotto da un team di tre persone (Angelo, Matteo,
Federico) nell'arco di pi\`u sessioni, da marzo a giugno 2026. Tutto il
materiale \`e stato tracciato in un repository versionato che raccoglie la
documentazione tecnica (\emph{Bill of Materials}, datasheet), i piani di
manutenzione, i log di volo in formato ULog di PX4 e gli strumenti di analisi
(script di plotting e una pipeline di visualizzazione 3D in Foxglove). Questa
relazione \`e l'estratto ragionato di quel materiale.

L'impostazione scelta \`e \textbf{tematico-cronologica}: i capitoli sono
organizzati per tema, e all'interno di ciascun tema si segue l'ordine
temporale degli eventi. Questa scelta evita di frammentare lo stesso argomento
su pi\`u giornate --- il sottosistema GPS, ad esempio, \`e stato toccato in
cinque sessioni diverse, ma viene trattato in modo unitario nel
Capitolo~\ref{cap:diagnostici}.

\section{Sintesi del sistema}

Il velivolo \`e un esacottero a sei rotori su frame DJI F550, con
distribuzione di potenza integrata nella piastra inferiore. Il cervello di
bordo \`e un flight controller Pixhawk 6X, dotato di tre IMU ridondanti
gestite in \emph{sensor voting} dal firmware PX4. La navigazione si appoggia a
un modulo GNSS montato su mast (per allontanarlo dalle interferenze
elettromagnetiche dei motori) con supporto a correzione RTK. La
configurazione, il monitoraggio e l'analisi delle missioni avvengono tramite
la ground station QGroundControl, collegata via telemetria a 433\,MHz.
Il Capitolo~\ref{cap:sistema} dettaglia l'architettura e ne riepiloga la
\emph{Bill of Materials}.

\figurasegnaposto{Esacottero DJI F550 in configurazione di volo (immagine di
apertura).}{Foto -- da scattare}{fig:drone-apertura}

\section{Struttura della relazione}

La relazione \`e organizzata come segue. Il Capitolo~\ref{cap:sistema}
descrive l'architettura hardware. Il Capitolo~\ref{cap:configurazione} ripercorre
la configurazione via QGroundControl e PX4, con le problematiche incontrate. Il
Capitolo~\ref{cap:acquisizione} --- cardine del progetto --- spiega \emph{cosa}
viene loggato e \emph{a che frequenza}. Il Capitolo~\ref{cap:campagna} presenta
la campagna di volo e tutti i tipi di prova eseguiti. Il
Capitolo~\ref{cap:diagnostici} raccoglie i casi diagnostici pi\`u significativi.
Il Capitolo~\ref{cap:osservazioni} propone alcune osservazioni sui dati
acquisiti. Il Capitolo~\ref{cap:foxglove} presenta, come contributo extra, lo
strumento di visualizzazione realizzato. Le conclusioni e le appendici chiudono
il documento.
